%===============================================================================
% 中英文摘要
%===============================================================================

% ---------- 中文摘要 ----------
\newpage
\thispagestyle{plain}
\begin{center}
    \heiti\fontsize{18pt}{27pt}\selectfont 摘\hskip1em 要
\end{center}
\vspace{2.8mm}

\fontsize{12pt}{18pt}\selectfont
随着嵌入式系统对图形显示、数据采集与缓存容量的需求不断增长，STM32 系列
微控制器内部 SRAM 已难以满足应用要求。可变静态存储器控制器（Flexible Static
Memory Controller, FSMC）作为 STM32F4 系列内置的并行存储器接口外设，
支持 SRAM、NOR Flash、PSRAM、NAND Flash 以及 LCD 等多种器件的连接，
其时序参数的合理配置直接决定了访问带宽与系统稳定性。

本文围绕 STM32F407 平台的 FSMC 外设展开研究，分析了 FSMC 内部结构与
四个 Bank 的地址映射机制，推导了 SRAM 模式下读写时序参数（ADDSET、
ADDHLD、DATAST、CLKDIV、BUSTURN）与系统总线频率的约束关系；提出
一种自适应时序参数搜索方法，在保证访问稳定性的前提下最小化访问周期；
基于 STM32F407 + IS62WV51216 SRAM + ILI9341 LCD 的硬件平台搭建了
测试系统，对不同时序配置下的读写带宽进行了实测。

实验结果表明，相比厂家推荐的保守时序配置，本文方法可将 SRAM 写带宽
从 35.2 MB/s 提升至 58.7 MB/s，提升约 67\%；LCD 帧刷新延迟从 28 ms
降至 18 ms。研究成果可为基于 STM32 的高速外设接口设计提供参考。

\vspace{\baselineskip}
\noindent{\heiti\fontsize{14pt}{21pt}\selectfont 关键词：}
{\fontsize{12pt}{18pt}\selectfont STM32；FSMC；外扩存储器；时序参数；接口设计}

% ---------- 英文摘要 ----------
\newpage
\thispagestyle{plain}
\begin{center}
    \bfseries\fontsize{18pt}{27pt}\selectfont Abstract
\end{center}
\vspace{2.8mm}

\fontsize{12pt}{18pt}\selectfont
With the increasing demand of embedded systems for graphical display, data
acquisition, and large-capacity caching, the on-chip SRAM of STM32
microcontrollers can no longer satisfy modern application requirements. The
Flexible Static Memory Controller (FSMC), an on-chip parallel memory
interface peripheral integrated in the STM32F4 series, supports a wide range
of external devices including SRAM, NOR Flash, PSRAM, NAND Flash, and
LCD modules. Proper configuration of its timing parameters directly
determines the achievable access bandwidth and system stability.

This thesis investigates the FSMC peripheral on the STM32F407 platform.
The internal structure of the FSMC and the address-mapping mechanism of its
four banks are analyzed, and the constraints between read/write timing
parameters (ADDSET, ADDHLD, DATAST, CLKDIV, BUSTURN) and the bus
frequency are derived for the SRAM mode. An adaptive timing-parameter
search method is proposed to minimize the access cycle while preserving
stability. A test platform consisting of an STM32F407, an IS62WV51216
SRAM, and an ILI9341 LCD is implemented to measure the bandwidth under
various timing configurations.

Experimental results show that, compared with the conservative timing
configuration recommended by the manufacturer, the proposed method
improves the SRAM write bandwidth from 35.2 MB/s to 58.7 MB/s
(approximately 67\%), and reduces the LCD frame-refresh latency from
28 ms to 18 ms. The results provide a useful reference for designing
high-speed peripheral interfaces on STM32 platforms.

\vspace{\baselineskip}
\noindent{\bfseries\fontsize{14pt}{21pt}\selectfont Key Words: }
{\fontsize{12pt}{18pt}\selectfont STM32; FSMC; External Memory; Timing Parameters; Interface Design}
